Dense per-voxel annotation is the binding constraint on CT segmentation.
Vision--language models trained on image--text pairs offer a route around it,
but their contrastive objective aligns pooled image and caption embeddings, so
individual patch tokens are never required to be spatially discriminative:
reading per-patch similarity off such a model recovers 0.003 mIoU. We make that
spatial property an explicit training objective, aligning each ViT patch token
with the organ's text embedding against a downsampled mask, so per-patch cosine
similarity becomes directly readable as a segmentation map at inference, with no
gradients, no classifier, and no second-stage mask refiner. On 11 organ classes
from TotalSegmentator this reaches 0.653 mIoU and 0.757 Dice using no per-pixel
labels at inference; a calibration correction contributes 0.059 of that with no
retraining. A prompt ablation shows the text is load-bearing rather than a class
index, since corrupting positional language collapses exactly the classes that
depend on laterality and leaves uniquely-positioned organs unaffected. The same
ablation shows elaboration is not where the value lies: a bare class name
recovers 96.9\% of a full multi-sentence anatomical description, and a one-line
generic caption matches it outright. An accidental two-word prompt error
revealed that failures are directional rather than degradative, producing a
confident, clean segmentation of the wrong organ, which is the more dangerous
failure mode because nothing in the output signals it. Carrying the frozen model
onto an unseen dataset and protocol, we compute a bone-metastasis proximity
index over Spine-Mets-CT-SEG: across 180 lesions in 46 patients the correct
organ ranks 1.78 on average against a chance rank of 2.91 ($p < 0.0001$,
patient-clustered permutation test), beating chance at every level from T1 to
T11 and concentrating its failures in two of eleven classes. We
report two pre-registered calibration variants that were run and not adopted,
each with a diagnosed mechanism, along with a ground-truth error in our own
evaluation table found and corrected from anatomical literature, and a
calibration leak found and repaired. Both corrections changed headline numbers.
